% ============================================================
%  ENCODING E FONT
% ============================================================
\usepackage{fontspec}
\setmainfont{TeX Gyre Pagella}
% ============================================================
%  LINGUE
% ============================================================
\usepackage{polyglossia}
\setmainlanguage{italian}
\setotherlanguage[variant=ancient]{greek}

% ============================================================
%  MATEMATICA
% ============================================================
\usepackage{amsmath}
\usepackage{amssymb}
\usepackage{amsfonts}
\usepackage{amsthm}
\usepackage{mathtools}
\usepackage{physics}
\usepackage{yhmath}
\usepackage{cancel}
\usepackage{bbm}

% unicode-math va caricato dopo amsmath/amssymb/amsfonts (altrimenti conflitti
% tipo "\eth already defined"); fornisce il font matematico per XeLaTeX.
\usepackage{unicode-math}
\setmathfont{TeX Gyre Pagella Math}

%dimensionimatematica
\makeatletter
\DeclareMathSizes{\@xpt}{10}{7.5}{6}    % testo 10pt
\DeclareMathSizes{\@xipt}{10}{7.5}{6}   % testo 11pt
\DeclareMathSizes{\@xiipt}{10}{7.5}{6}  % testo 12pt
\makeatother

% Spaziatura display math
\setlength{\abovedisplayskip}{4pt}
\setlength{\belowdisplayskip}{4pt}
\setlength{\abovedisplayshortskip}{2pt}
\setlength{\belowdisplayshortskip}{2pt}

% Comandi matematici
\newcommand{\R}{\mathbb{R}}
\newcommand{\sse}{ \ \textup{sse} \ }
\newcommand{\tc}{ \textup{ t.c. }}
\newcommand{\under}[1]{_{\textup{#1}}}
\newcommand{\sopra}[1]{^{\textup{#1}}}
\newcommand{\tonde}[1]{\left(#1\right)}
\newcommand{\quadre}[1]{\left[#1\right]}
\newcommand{\graffe}[1]{\left\lbrace #1 \right\rbrace}
\newcommand{\et}{\wedge}
\newcommand{\vel}{\vee}
\newcommand{\aut}{\dot{\vee}}
\newcommand{\quindi}{\quad \Rightarrow \quad}
\newcommand{\ubar}[1]{\underline{\mathord{#1}}}

% ============================================================
%  UNITÀ DI MISURA
% ============================================================
\usepackage{siunitx}
\sisetup{
  detect-all,
  output-decimal-marker = {,}
}
\DeclareSIPrefix\micro{μ}{-6}
% ============================================================
%  LAYOUT E GEOMETRIA
% ============================================================
\usepackage[centering, includeheadfoot]{geometry}

\newlength{\margineEsterno}
\newlength{\margineInterno}
\newlength{\margineAlto}
\newlength{\margineBasso}

\setlength{\margineEsterno}{3.5cm}
\setlength{\margineInterno}{1cm}
\setlength{\margineAlto}{1cm}
\setlength{\margineBasso}{1cm}

\geometry{
  paperwidth  = 170mm,
  paperheight = 240mm,
  top         = \margineAlto,
  bottom      = \margineBasso,
  inner       = \margineInterno,
  outer       = \margineEsterno,
  twoside,
}

% Dimensioni margine per QR e note
\newlength{\maxQR}
\newlength{\sizeQR}
\newlength{\largQR}

\setlength{\maxQR}{2.5cm}
\setlength{\marginparsep}{3mm}
\setlength{\marginparwidth}{\dimexpr\margineEsterno - \marginparsep\relax}
\setlength{\largQR}{\dimexpr\marginparwidth - 6mm\relax}

\newcommand{\calcsizeQR}{%
  \setlength{\sizeQR}{\dimexpr\marginparwidth - 10pt\relax}%
  \ifdim\sizeQR>\maxQR \setlength{\sizeQR}{\maxQR}\fi
}

% Spaziatura testo
\usepackage{setspace}
\setstretch{1.18}
\usepackage{parskip}
\setlength{\parindent}{0pt}
\setlength{\parskip}{0.25\baselineskip plus 1pt minus 1pt}

%Sforare le immagini nel margine esterno
\newcommand{\sforaEsterno}[1]{%
  \checkoddpage
  \ifoddpage
    \begin{adjustwidth}{0pt}{-\margineEsterno}#1\end{adjustwidth}%
  \else
    \begin{adjustwidth}{-\margineEsterno}{0pt}#1\end{adjustwidth}%
  \fi
}

\raggedbottom

% ============================================================
%  INTESTAZIONI E PIÈ DI PAGINA
% ============================================================
\usepackage{fancyhdr}
\pagestyle{fancy}
\fancyhf{}
\fancyfoot[C]{\thepage}
\renewcommand{\headrulewidth}{0pt}
\renewcommand{\footrulewidth}{0pt}
\let\cleardoublepage\clearpage
\renewcommand{\cleardoublepage}{\clearpage}

% ============================================================
%  TITOLI E SEZIONI
% ============================================================
\usepackage{titlesec}
\usepackage{tocloft}

% Part
\renewcommand{\part}[1]{%
  \refstepcounter{part}
  \clearpage
  \thispagestyle{empty}
  \begin{center}
    \Huge\bfseries Part \thepart \\[1em] #1
  \end{center}
  \vspace{2em}
  \addcontentsline{toc}{part}{Part \thepart: #1}
}

% Helper: salva il nome del capitolo E lo stampa
\newcommand{\savechaptertitle}[1]{\gdef\currentchaptername{#1}#1}

\titleformat{\chapter}[block]
  {\normalfont\LARGE\bfseries\color{cyan}}
  {\rule{\textwidth}{.5pt}\\[0.1em]\centering\Roman{chapter}.\,}
  {1em}
  {\huge\scshape\centering\savechaptertitle}
  [\vspace{-0.2em}\rule{\textwidth}{0.5pt}]

% ============================================================
%  CAPITOLO CON IMMAGINE OPZIONALE
% ============================================================
\newcommand{\chapterimage}[3]{%
  \clearpage
  \thispagestyle{fancy}%
  \begin{adjustwidth}{-\margineEsterno}{0pt}%
    \includegraphics[width=\dimexpr\textwidth+\margineEsterno\relax]{#1}%
    \par\vspace{2pt}%
    {\small\itshape\color{gray} #2}%
  \end{adjustwidth}%
  \vspace{0.5em}%
  \refstepcounter{chapter}%
  \gdef\currentchaptername{#3}%
  \addcontentsline{toc}{chapter}{\protect\numberline{\thechapter}#3}%
  \begin{center}
    {\color{cyan}\rule{\textwidth}{.5pt}}\\[0.2em]
    {\color{cyan}\LARGE\bfseries\Roman{chapter}.\,}\\[0.1em]
    {\color{cyan}\huge\scshape #3}\\[-0.3em]
    {\color{cyan}\rule{\textwidth}{.5pt}}
  \end{center}
  \vspace{0.5em}%
}

\titleformat{\section}[block]
  {\normalfont\Large\bfseries\color{cyan}}{\thesection}{1em}{}

\titleformat{\subsection}[block]
  {\normalfont\large\bfseries\color{cyan}}{\thesubsection}{1em}{}

\titleformat{\subsubsection}[block]
  {\normalfont\bfseries\color{cyan}}{\thesubsubsection}{1em}{}

\titleformat{\paragraph}[block]
  {\normalfont\bfseries\color{cyan}}{}{0em}{}

\titlespacing*{\chapter}   {0pt}{-30pt}{8pt}
\titlespacing*{\section}   {0pt}{0.8\baselineskip}{0.3\baselineskip}
\titlespacing*{\subsection}{0pt}{0.6\baselineskip}{0.2\baselineskip}
\setlength{\cftbeforechapskip}{4pt}

% ============================================================
%  LISTE
% ============================================================
\usepackage[inline]{enumitem}
\setlist[itemize,1]{label=\textcolor{gray!70}{\rule{0.8ex}{0.8ex}}}
\setlist{itemsep=0pt, parsep=0pt, topsep=3pt, partopsep=0pt}
\setlength{\topsep}{2pt}
\newlist{myenum}{enumerate*}{1}
\setlist[myenum]{label=\arabic*, itemjoin={\hspace{1.5cm}}}
\newlist{sceltamultipla}{enumerate}{1}
\setlist[sceltamultipla,1]{%
  label=\textcolor{gray}{\fbox{\textbf{\footnotesize\makebox[1.2em]{\Alph*}}}},
  leftmargin=*,
  itemsep=1pt,
  topsep=2pt
}

% ============================================================
%  COLORI
% ============================================================
\usepackage{xcolor}

\definecolor{ciano}           {RGB}{0,   180, 200}
\definecolor{sfondoImportante}{RGB}{255, 252, 220}
\definecolor{coloreAnimazione}{RGB}{0,   150, 136}
\definecolor{sfondoAnimazione}{RGB}{225, 245, 243}
\definecolor{gold}            {RGB}{255, 215,   0}
\definecolor{copper}          {RGB}{104, 115,  51}

\definecolor{buildSub}  {HTML}{FFF0F5}  \definecolor{buildTitle} {HTML}{C71585}
\definecolor{readSub}   {HTML}{F2F2F2}  \definecolor{readTitle}  {HTML}{333333}
\definecolor{deriveSub} {HTML}{FFF9E6}  \definecolor{deriveTitle}{HTML}{B8860B}
\definecolor{exSub}     {HTML}{F0F7FF}  \definecolor{exTitle}    {HTML}{005C99}
\definecolor{impSub}    {HTML}{FFF0F0}  \definecolor{impTitle}   {HTML}{CC0000}
\definecolor{warnSub}   {HTML}{FFF0F0}  \definecolor{warnTitle}  {HTML}{B22222}

\definecolor{sfondoRiassunto}{RGB}{215, 242, 240}
\definecolor{intestRiassunto}{RGB}{  0, 145, 132}
\definecolor{rsezColor}      {RGB}{  0, 160, 146}

% ============================================================
%  TIKZ E GRAFICI
% ============================================================
\usepackage{tikz}
\usepackage{tikz-3dplot}
\usepackage{pgfplots}
\usepackage{circuitikz}
\usepackage{tikzsymbols, fontawesome5}

\usetikzlibrary{
  decorations.pathreplacing, decorations.markings, decorations.pathmorphing,
  angles, quotes, arrows.meta, calc, positioning,
  shapes, shapes.symbols, shadows, fit,
  backgrounds, matrix, patterns, tikzmark
}

\pgfplotsset{compat=1.17}
\usepgfplotslibrary{colormaps}
\tdplotsetmaincoords{70}{110}

% Simbolo caution
\newcommand{\caution}[1]{%
  \begingroup
  \pgfmathsetmacro{\s}{#1}%
  \pgfmathsetmacro{\Rr}{1.6*\s}%
  \pgfmathsetmacro{\Lwd}{2*\s}%
  \pgfmathsetmacro{\FSpt}{40*\s}%
  \pgfmathsetmacro{\BLSpt}{1.2*\FSpt}%
  \tikz[baseline=(bang.base)]{
    \draw[line width=\Lwd mm, red, fill=yellow!30]
      (90:\Rr) -- (210:\Rr) -- (330:\Rr) -- cycle;
    \node (bang) at (0,0)
      {\fontsize{\FSpt pt}{\BLSpt pt}\selectfont\bfseries !};
  }%
  \endgroup
}

% ============================================================
%  BOX: TCOLORBOX
% ============================================================
\usepackage[many]{tcolorbox}

\tcbset{
  mystyle/.style={
    enhanced, breakable,
    boxrule=0pt, frame hidden, arc=2mm,
    top=4mm, bottom=4mm, left=4mm, right=4mm,
    fonttitle=\bfseries\sffamily, coltitle=white,
    attach boxed title to top left={yshift=-2mm, xshift=0mm},
    boxed title style={
      empty,
      top=1.5mm, bottom=1.5mm, left=3mm, right=8mm,
      underlay={
        \fill[#1, rounded corners=2mm]
          (frame.south west) -- (frame.north west) --
          ([xshift=-5mm]frame.north east) -- (frame.south east) -- cycle;
      }
    }
  }
}

\newtcolorbox{mdwarning}[1][]{
  mystyle={warnTitle}, colback=warnSub,
  title={\ifstrempty{#1}{Attenzione}{#1}}
}
\newtcolorbox{mdbuild}[1][]{
  mystyle={buildTitle}, colback=buildSub,
  title={\ifstrempty{#1}{Costruire la formula}{#1}}
}
\newtcolorbox{mdread}[1][]{
  mystyle={readTitle}, colback=readSub,
  title={\ifstrempty{#1}{Leggere la formula}{#1}}
}
\newtcolorbox{mdderive}[1][]{
  mystyle={deriveTitle}, colback=deriveSub,
  title={\ifstrempty{#1}{Derivare la formula}{#1}}
}
\newtcolorbox{mdexample}[1][]{
  mystyle={exTitle}, colback=exSub,
  title={\ifstrempty{#1}{Esempio guidato}{#1}}
}

\newtcolorbox{mdimportant}[2][]{%
  enhanced, breakable,
  boxrule=0pt, frame hidden,
  colback=sfondoImportante, arc=1mm,
  top=4mm, bottom=4mm, left=6mm, right=4mm,
  overlay={
    \checkoddpage
    \ifoddpage
      \draw[ciano, line width=2pt] (frame.south east)--(frame.north east);
      \draw[ciano, line width=2pt] (frame.east)--([xshift=5mm]frame.east);
      \draw[ciano, fill=white, line width=2pt] ([xshift=5mm]frame.east) circle (4pt);
      \node[anchor=west, text=ciano, font=\bfseries\sffamily\small,
            align=left, text width=\dimexpr\largQR-2mm\relax]
        at ([xshift=7mm]frame.east) {#2};
    \else
      \draw[ciano, line width=2pt] (frame.south west)--(frame.north west);
      \draw[ciano, line width=2pt] ([xshift=-5mm]frame.west)--(frame.west);
      \draw[ciano, fill=white, line width=2pt] ([xshift=-5mm]frame.west) circle (4pt);
      \node[anchor=east, text=ciano, font=\bfseries\sffamily\small,
            align=right, text width=\dimexpr\largQR-2mm\relax]
        at ([xshift=-7mm]frame.west) {#2};
    \fi
  },
  #1
}

\newsavebox{\mdanimationQRbox}

\newtcolorbox{mdanimation}[2][Animazione]{%
  enhanced, breakable,
  boxrule=0pt, frame hidden,
  colback=sfondoAnimazione, arc=2mm,
  top=4mm, bottom=4mm, left=4mm, right=4mm,
  fonttitle=\bfseries\sffamily,
  coltitle=white,
  title={#1},
  attach boxed title to top left={yshift=-2mm, xshift=0mm},
  boxed title style={
    empty,
    top=1.5mm, bottom=1.5mm, left=3mm, right=8mm,
    underlay={
      \fill[coloreAnimazione, rounded corners=2mm]
        (frame.south west) -- (frame.north west) --
        ([xshift=-5mm]frame.north east) -- (frame.south east) -- cycle;
    }
  },
  overlay={
    \ifx\relax#2\relax\else
      \checkoddpage
      \calcsizeQR
      \ifoddpage
        \fill[sfondoAnimazione]
          ([xshift=-1mm, yshift=-2mm]frame.north east) --
          ([xshift=3mm, yshift=1mm]frame.north east) --
          ([xshift=3mm, yshift=-\the\sizeQR-0pt]frame.north east) --
          ([xshift=-1mm, yshift=-\the\sizeQR+3mm]frame.north east) -- cycle;
        \node[anchor=north west, draw=coloreAnimazione, line width=1.5pt,
              fill=sfondoAnimazione, inner sep=4pt, rounded corners=1mm]
          at ([xshift=3mm, yshift=2mm]frame.north east) {%
            \begin{minipage}{\the\largQR}\centering
              \qrcode[height=\the\sizeQR]{#2}
            \end{minipage}%
          };
      \else
        \fill[sfondoAnimazione]
          ([xshift=1mm, yshift=-2mm]frame.north west) --
          ([xshift=-3mm, yshift=1mm]frame.north west) --
          ([xshift=-3mm, yshift=-\the\sizeQR-0pt]frame.north west) --
          ([xshift=1mm, yshift=-\the\sizeQR+3mm]frame.north west) -- cycle;
        \node[anchor=north east, draw=coloreAnimazione, line width=1.5pt,
              fill=sfondoAnimazione, inner sep=4pt, rounded corners=1mm]
          at ([xshift=-3mm, yshift=2mm]frame.north west) {%
            \begin{minipage}{\the\largQR}\centering
              \qrcode[height=\the\sizeQR]{#2}
            \end{minipage}%
          };
      \fi
    \fi
  },
}

% ============================================================
%  BOX: MDFRAMED
% ============================================================
\usepackage{mdframed}

\newenvironment{cyanmargin}[1][]
  {\begin{mdframed}[topline=false,bottomline=false,rightline=false,
    linecolor=cyan,linewidth=.05cm,frametitlefont=\itshape]#1\ignorespaces}
  {\end{mdframed}}

\newenvironment{mdexercise}[1][]
  {\begin{mdframed}[topline=false,bottomline=false,rightline=false,
    linecolor=cyan,linewidth=.05cm,frametitlefont=\itshape]
   \textbf{Esercizio. }#1\ignorespaces}
  {\end{mdframed}}

\newenvironment{mdexam}[1][]
  {\begin{mdframed}[topline=false,bottomline=false,rightline=false,
    linecolor=magenta,linewidth=2,frametitlefont=\itshape]
   \textbf{Allenamento per la verifica. }#1\ignorespaces}
  {\end{mdframed}}

\newenvironment{mdproof}[1][]
  {\begin{mdframed}[topline=false,bottomline=false,rightline=false,
    linecolor=black,frametitle=Proof.,frametitlefont=\itshape]#1\ignorespaces}
  {\unskip\begin{flushright}\qedsymbol\end{flushright}\end{mdframed}}

\newenvironment{nestedproof}[1][\small]
  {\begin{mdframed}[topline=false,bottomline=false,rightline=false,
    linecolor=black,frametitle=Proof.,frametitlefont=\itshape]#1\ignorespaces}
  {\unskip\begin{flushright}\qedsymbol\end{flushright}\end{mdframed}}

% ============================================================
%  TEOREMI
% ============================================================
\newtheoremstyle{flemma}
  {3pt}{3pt}
  {\itshape\fontfamily{qpl}\selectfont}
  {1pt}{\bfseries}{ \space}{ }{}

\theoremstyle{definition}
\newtheorem*{definition}  {Definizione}
\newtheorem*{proposition} {Proposizione}
\newtheorem*{example}     {Esempio}
\newtheorem*{remark}      {Osservazione}
\newtheorem*{nota}        {Nota}
\newtheorem*{application} {Applicazione}

\theoremstyle{flemma}
\newtheorem*{property}  {Proprietà}
\newtheorem*{law}       {Legge}
\newtheorem*{lemma}     {Lemma}
\newtheorem*{theorem}   {Teorema}
\newtheorem*{corollary} {Corollario}
\newtheorem*{principle} {Principio}

\renewcommand\qedsymbol{\fontfamily{qpl}\selectfont\emph{Q.E.D.}}

% ============================================================
%  ESERCIZI E SOLUZIONI (ambienti classici)
% ============================================================
\usepackage{needspace}
\usepackage{xparse}
\usepackage{pifont}
\usepackage{environ}

\newcounter{exercise}
\renewcommand{\theexercise}{\arabic{exercise}}

% Stelline — nere, senza fbox
\newcommand{\onestar}  {\ding{72}\ding{73}\ding{73}\space}
\newcommand{\twostars} {\ding{72}\ding{72}\ding{73}\space}
\newcommand{\threestars}{\ding{72}\ding{72}\ding{72}\space}

\newcommand{\setexdifficulty}[1]{%
  \IfNoValueTF{#1}
    {\def\exdifficulty{\twostars}}
    {\def\exdifficulty{#1}}%
}

\NewDocumentEnvironment{exercise}{ o }{%
  \refstepcounter{exercise}\par\medskip
  \setexdifficulty{#1}\noindent
  \hypertarget{ex:\theexercise}{}\label{ex:\theexercise}%
  \textbf{Esercizio \theexercise\ }\exdifficulty\quad
  \hyperlink{sol:\theexercise}{\textit{(vai alla soluzione)}}\quad\ignorespaces
}{\par\medskip}

\NewDocumentEnvironment{question}{ o }{%
  \refstepcounter{exercise}\par\medskip
  \setexdifficulty{#1}\noindent
  \hypertarget{ex:\theexercise}{}\label{ex:\theexercise}%
  \textbf{Domanda di Teoria \theexercise\ }\exdifficulty\quad
  \hyperlink{sol:\theexercise}{\textit{(vai alla soluzione)}}\quad\ignorespaces
}{\par\medskip}

\NewDocumentEnvironment{exercisenosol}{ o }{%
  \refstepcounter{exercise}\par\medskip
  \setexdifficulty{#1}\noindent
  \hypertarget{ex:\theexercise}{}\label{ex:\theexercise}%
  \textbf{Esercizio \theexercise\ }\exdifficulty\quad\ignorespaces
}{\par\medskip}

\NewDocumentEnvironment{questionnosol}{ o }{%
  \refstepcounter{exercise}\par\medskip
  \setexdifficulty{#1}\noindent
  \hypertarget{ex:\theexercise}{}\label{ex:\theexercise}%
  \textbf{Domanda di teoria \theexercise\ }\exdifficulty\quad\ignorespaces
}{\par\medskip}

\NewDocumentEnvironment{expsolution}{ o }{%
  \Needspace{6\baselineskip}%
  \IfNoValueTF{#1}{\edef\resnum{\theexercise}}{\edef\resnum{#1}}%
  \par\bigskip
  \hypertarget{sol:\resnum}{}\label{sol:\resnum}%
  \noindent\textbf{Soluzione \hyperlink{ex:\resnum}{\resnum}}\quad
  \hyperlink{ex:\resnum}{\textit{(torna all'esercizio)}}\par\medskip
}{\par\medskip}

% solution: comportamento duale (inline grigio fuori da rigaesercizi,
%           accumulatore dentro rigaesercizi — vedi sotto)
\newenvironment{solution}
  {\par\noindent\hfill\small\color{gray}}
  {\par}

% ============================================================
%  MODALITÀ ESERCIZI (due colonne + testata ciano)
% ============================================================
\usepackage{multicol}
\usepackage{colortbl}

\newcommand{\exheadrule}[2]{%
  \setlength{\headheight}{20pt}%
  \fancyhead[RE]{%
    \colorbox{ciano}{%
      \parbox[c][14pt][c]{\textwidth}{%
        \color{white}\bfseries\sffamily\small #1\hfill}}}%
  \fancyhead[LO]{%
    \colorbox{ciano}{%
      \parbox[c][14pt][c]{\textwidth}{%
        \color{white}\bfseries\sffamily\small\hfill #2}}}%
  \fancyhead[RO,LE]{}%
  \renewcommand{\headrulewidth}{0pt}%
}

\newif\ifinsidemodalita
\insidemodalitafalse

\newenvironment{modalitaesercizi}{%
  \insidemodalitatrue
  \edef\exChapNum{\thechapter}%
  \edef\exChapName{\currentchaptername}%
  \renewcommand{\Needspace}[1]{}%
  \newgeometry{
    top=0cm, bottom=.8cm,
    inner=1.2cm, outer=1.2cm,
    includeheadfoot,
  }%
  \pagestyle{fancy}\fancyhf{}\fancyfoot[C]{\thepage}%
  \exheadrule%
    {CAPITOLO \exChapNum \enspace\textbar\enspace \exChapName}%
    {CAPITOLO \exChapNum \enspace\textbar\enspace ESERCIZI}%
  \setlength{\columnsep}{1cm}%
  \raggedcolumns
  \begin{multicols}{2}%
}{%
  \end{multicols}%
  \restoregeometry%
  \pagestyle{fancy}\fancyhf{}\fancyhead{}\fancyfoot[C]{\thepage}%
  \renewcommand{\headrulewidth}{0pt}%
  \insidemodalitafalse
}

\newcommand{\pausamulticol}{%
  \ifinsidemodalita\end{multicols}\par\vspace{2pt}\fi
}
\newcommand{\riprendimulticol}{%
  \ifinsidemodalita\par\vspace{2pt}\begin{multicols}{2}\fi
}

% ============================================================
%  AMBIENTE truthorfalse
%  Da aggiungere nel preambolo, dopo il caricamento di
%  tcolorbox (già presente) e enumitem (già presente).
%
%  Uso:
%    \begin{truthorfalse}
%      \tfitem{testo dell'affermazione}
%    \end{truthorfalse}
%
%  Le caselle V / F vengono inserite automaticamente
%  a destra di ogni voce.
% ============================================================

% Casellina grigia piccola (V oppure F).
% Si usa \newtcbox invece di \tcbox[...]{...} inline:
% \tcbox con opzioni esplicite dentro \newcommand causa
% conflitti nel parsing degli argomenti di tcolorbox.
% \newtcbox pre-registra le opzioni ed è privo di ambiguità.
\newtcbox{\tfvf}{%
  on line,
  arc=0.8pt, outer arc=0.8pt,
  boxrule=0.3pt, boxsep=0pt,
  left=2pt, right=2pt, top=0.8pt, bottom=0.8pt,
  colframe=gray!40, colback=gray!10,
  fontupper={\fontsize{6}{6}\selectfont\bfseries\color{gray!65}}%
}

% Singola voce: testo → caselle V F allineate a destra
\newcommand{\tfitem}[1]{%
  \item #1\nobreak\hfill\tfvf{V}\kern2pt\tfvf{F}%
}

% Lista vera/falsa
\newenvironment{truthorfalse}{%
  \begin{enumerate}[%
    label  = \textbf{\alph*.},%
    leftmargin = *,%
    itemsep = 3pt,%
    topsep  = 2pt,%
    parsep  = 0pt%
  ]%
}{%
  \end{enumerate}%
}
% ============================================================
%  Flag per raccolta soluzioni
% ============================================================
% ============================================================
%  *** SOSTITUIRE TUTTO da qui fino a fine \newenvironment{exrcell} ***
% ============================================================

% Flag per raccolta soluzioni
\newif\ifincollectmode  \incollectmodefalse
\newif\iffirstsolution  \firstsolutiontrue

% ============================================================
%  CORREZIONE BUG soluzioni duplicate con rigaesercizi{3}
%
%  Causa: il vecchio meccanismo con \newtoks memorizzava
%  \the\exrsoltemp come riferimento LAZY dentro \exrsoltoks.
%  Quando la terza soluzione sovrascriveva \exrsoltemp,
%  entrambi i riferimenti espandevano lo stesso valore:
%  risultato [sol1, sol3, sol3] invece di [sol1, sol2, sol3].
%
%  Fix: ogni soluzione viene copiata permanentemente in
%  \exrsol@1, \exrsol@2, \exrsol@3  via \let, che fotografa
%  \BODY al momento della chiamata.  Ridefinizioni successive
%  di \BODY non toccano le copie già salvate.
% ============================================================

\newcounter{exrsolcount}   % conta le soluzioni della riga corrente

\RenewEnviron{solution}{%
  \ifincollectmode
    \iffirstsolution
      \global\firstsolutionfalse
    \fi
    \stepcounter{exrsolcount}%
    % Copia permanente: \exrsol@N ottiene il significato corrente
    % di \BODY; eventuali ridefinizioni future di \BODY non la toccano
    \expandafter\global\expandafter\let
      \csname exrsol@\theexrsolcount\endcsname\BODY
  \else
    \par\noindent\hfill\small\color{gray}\BODY\par
  \fi
}

% ============================================================
%  Dimensioni e box ausiliari (invariati)
% ============================================================
\newsavebox{\exrfulltmpbox}
\newsavebox{\exrsolbox}
\newsavebox{\exrstarbox}
\newlength{\exrlabelwidth}
\newlength{\exrcellwidth}
\newlength{\exrsolcolwidth} \setlength{\exrsolcolwidth}{3.5cm}
\newlength{\exrcellsep}     \setlength{\exrcellsep}{6mm}
\newlength{\exrlabelgap}    \setlength{\exrlabelgap}{4mm}
\newlength{\exrtotalgap}
\newcounter{exrowcell}
\newcounter{exrowmaxcolsctr}

% ============================================================
%  Stampa le soluzioni accumulate come  [sol1; sol2; sol3]
%  Gestisce fino a 3 soluzioni per riga (massimo supportato).
% ============================================================
\newcommand{\exrprintsol}{%
  \savebox{\exrsolbox}{%
    \color{gray!70!black}%
    \ifnum\value{exrsolcount}>0\relax
      \csname exrsol@1\endcsname
      \ifnum\value{exrsolcount}>1\relax
        ;\space\csname exrsol@2\endcsname
        \ifnum\value{exrsolcount}>2\relax
          ;\space\csname exrsol@3\endcsname
        \fi
      \fi
    \fi
  }%
  \color{gray!70!black}%
  $\left[\rule[-\dp\exrsolbox]{0pt}%
    {\dimexpr\ht\exrsolbox+\dp\exrsolbox\relax}\right.$%
  \usebox{\exrsolbox}%
  $\left.\rule[-\dp\exrsolbox]{0pt}%
    {\dimexpr\ht\exrsolbox+\dp\exrsolbox\relax}\right]$%
}

\newenvironment{rigaesercizi}[2][\twostars]{%
  \setcounter{exrowcell}{0}%
  \setcounter{exrowmaxcolsctr}{#2}%
  \setcounter{exrsolcount}{0}%
  \global\firstsolutiontrue%
  \incollectmodetrue%
  \refstepcounter{exercise}%
  \setexdifficulty{#1}%
  \hypertarget{ex:\theexercise}{}\label{ex:\theexercise}%
  \sbox{\exrstarbox}{{\scriptsize\exdifficulty}}%
  \setlength{\exrlabelwidth}{\dimexpr\wd\exrstarbox + 6pt\relax}%
  \setlength{\exrtotalgap}{#2\exrcellsep}%
  \advance\exrtotalgap by -\exrcellsep%
  \ifnum#2=1\relax
    \setlength{\exrcellwidth}{\dimexpr
      \linewidth - \exrlabelwidth - \exrlabelgap\relax}%
  \else
    \setlength{\exrcellwidth}{\dimexpr
      ( \linewidth
        - \exrlabelwidth
        - \exrlabelgap
        - \exrsolcolwidth
        - \exrcellsep
        - \exrtotalgap
      ) / #2\relax}%
  \fi
  \par\nopagebreak\vspace{2pt}%
  % \vbox\bgroup\noindent\leavevmode
  \vbox\bgroup\noindent\leavevmode\parfillskip=0pt
  {\setlength{\fboxsep}{0pt}%
   \vtop{\offinterlineskip
     \hbox{\colorbox{intestRiassunto}{%
       \color{white}\bfseries\small
       \makebox[\exrlabelwidth][c]{\rule{0pt}{11pt}\theexercise}}}%
     \hbox{\makebox[\exrlabelwidth][c]{\scriptsize\exdifficulty}}%
   }}%
  \hspace{\exrlabelgap}\small\raggedright\parfillskip=0pt
}{%
  \iffirstsolution\else
    \unskip\nobreak\hfill\allowbreak\null\hfill\mbox{\exrprintsol}%
  \fi
  \par\egroup
  \vspace{4pt}%
  \incollectmodefalse%
}

\newenvironment{exrcell}{%
  \stepcounter{exrowcell}%
  \ifnum\value{exrowmaxcolsctr}=1\relax
    \ignorespaces
  \else
    \vtop\bgroup
      \hsize=\exrcellwidth
      \binoppenalty=10000
      \relpenalty=10000
      \small\raggedright\ignorespaces
  \fi
}{%
  \ifnum\value{exrowmaxcolsctr}=1\relax
    \unskip
  \else
    \egroup
    \ifnum\value{exrowcell}<\value{exrowmaxcolsctr}\relax
      \hspace{\exrcellsep}%
    \fi
  \fi
}

% ============================================================
%  *** FINE sezione da sostituire ***
% ============================================================
% ============================================================
%  SEZIONI INTERNE AL RIASSUNTO
% ============================================================
\newcommand{\rsection}[1]{%
  \par\vspace{0.55\baselineskip}%
  \begin{tcolorbox}[
    enhanced, breakable,
    colback=rsezColor, colframe=rsezColor,
    arc=2mm, boxrule=0pt,
    top=2mm, bottom=2mm, left=5mm, right=5mm,
  ]
    \color{white}\bfseries\sffamily\large #1%
  \end{tcolorbox}%
  \vspace{0.1\baselineskip}%
}

\newcommand{\rsubsection}[1]{%
  \par\vspace{0.35\baselineskip}%
  \noindent{\color{intestRiassunto}\bfseries\sffamily #1}%
}

\newcommand{\rkeyword}[1]{%
  \textcolor{intestRiassunto}{\bfseries #1}%
}

% ============================================================
%  AMBIENTE RIASSUNTO
% ============================================================
\newenvironment{riassunto}{%
  \clearpage
  \pagecolor{sfondoRiassunto}%
  \newgeometry{
    top    = 1.8cm,
    bottom = 1.8cm,
    inner  = 1.5cm,
    outer  = 1.5cm,
    includeheadfoot,
  }%
  \pagestyle{fancy}\fancyhf{}%
  \setlength{\headheight}{22pt}%
  \fancyhead[LE,RO]{%
    \colorbox{intestRiassunto}{%
      \parbox[c][16pt][c]{\textwidth}{%
        \color{white}\bfseries\sffamily\small
        \hspace{3mm}CAPITOLO~\thechapter
        \enspace\textbar\enspace
        \MakeUppercase{\currentchaptername}%
        \hfill RIASSUNTO\hspace{3mm}%
      }%
    }%
  }%
  \fancyhead[LO,RE]{}%
  \renewcommand{\headrulewidth}{0pt}%
  \fancyfoot[C]{\color{intestRiassunto}\thepage}%
  \par\vspace{0.6em}%
  \begin{center}
    {\color{intestRiassunto!50!black}\rule{0.18\textwidth}{1pt}}%
    \quad
    {\color{intestRiassunto}\fontsize{27}{31}\selectfont\bfseries\scshape Riassunto}%
    \quad
    {\color{intestRiassunto!50!black}\rule{0.18\textwidth}{1pt}}\\[0.35em]
    {\color{rsezColor!70!white}\fontsize{9}{11}\selectfont\sffamily
      Capitolo~\thechapter\ \textbullet\ \currentchaptername}%
  \end{center}%
  \vspace{0.7em}%
}{%
  \clearpage
  \pagecolor{white}%
  \restoregeometry%
  \pagestyle{fancy}%
  \fancyhf{}\fancyhead{}\fancyfoot[C]{\thepage}%
  \renewcommand{\headrulewidth}{0pt}%
}

% ============================================================
%  AMBIENTE VERIFICA SUBITO
% ============================================================
\newenvironment{verificasubito}[2][Verifica subito!]{%
  \def\vsncols{#2}%
  \renewcommand{\Needspace}[1]{}%
  \checkoddpage%
  \ifoddpage%
    \def\sforoSin{0cm}%
    \def\sforoDes{0cm}%
  \else%
    \def\sforoSin{0cm}%
    \def\sforoDes{0cm}%
  \fi%
  \begin{tcolorbox}[
    enhanced, breakable,
    colback=white, colframe=ciano,
    boxrule=1.5pt, arc=3mm,
    top=4mm, bottom=4mm, left=4mm, right=4mm,
    enlarge left by=\sforoSin,
    enlarge right by=\sforoDes,
    title={\large\bfseries\sffamily #1},
    coltitle=white,
    attach boxed title to top left={yshift=-2mm, xshift=4mm},
    boxed title style={
      empty,
      top=1.5mm, bottom=1.5mm, left=3mm, right=8mm,
      underlay={
        \fill[ciano, rounded corners=2mm]
          (frame.south west) -- (frame.north west) --
          ([xshift=-5mm]frame.north east) -- (frame.south east) -- cycle;
      }
    },
  ]%
  \setlength{\columnsep}{1cm}%
  \raggedcolumns%
  \ifnum\vsncols>1\relax
    \begin{multicols}{\vsncols}%
  \fi
}{%
  \ifnum\vsncols>1\relax
    \end{multicols}%
  \fi
  \end{tcolorbox}%
}

% ============================================================
%  QR NEL MARGINE
% ============================================================
\usepackage{qrcode}
\usepackage{marginnote}
\usepackage{marginfix}

\newcommand{\marginqr}[2][]{%
  \par\Needspace{6\baselineskip}\calcsizeQR
  \marginpar{\centering
    \begin{tcolorbox}[enhanced, boxrule=1.5pt, colframe=ciano, colback=white,
      arc=1mm, top=4pt, bottom=4pt, left=3pt, right=3pt,
      boxsep=0pt, width=\marginparwidth]
      \if\relax\detokenize{#1}\relax\else
        {\footnotesize\setlength{\parskip}{0pt}\setstretch{1.0}%
         \raggedright\sffamily #1\par\vspace{3pt}}
      \fi
      \centering\qrcode[height=\sizeQR]{#2}%
    \end{tcolorbox}}%
}


% ============================================================
%  AMBIENTE PROVATU — box nel margine, colori riassunto
% ============================================================
\newsavebox{\provatubox}
\newenvironment{provatu}[1][Prova tu!]{%
  \begin{lrbox}{\provatubox}%
    \begin{minipage}{\marginparwidth}%
      \begin{mdframed}[
        skipabove=0pt, skipbelow=0pt,
        linecolor=intestRiassunto,
        linewidth=1.2pt,
        backgroundcolor=sfondoRiassunto,
        frametitle={\color{white}\small\bfseries\sffamily\hspace{2pt}#1},
        roundcorner=4pt,
        frametitlebackgroundcolor=intestRiassunto,
        frametitlerule=false,
        frametitleaboveskip=2pt,
        frametitlebelowskip=2pt,
        innertopmargin=4pt,
        innerbottommargin=4pt,
        innerleftmargin=4pt,
        innerrightmargin=4pt,
      ]%
      \small\raggedright\ignorespaces
}{%
      \end{mdframed}%
    \end{minipage}%
  \end{lrbox}%
  \marginnote{\usebox{\provatubox}}%
  \ignorespaces
}{\ignorespacesafterend}

% ============================================================
%  TESTO E TIPOGRAFIA
% ============================================================
\usepackage{csquotes}
\usepackage[normalem]{ulem}
\usepackage{soul}
\usepackage{tipa}

\setulcolor{gray}
\renewcommand{\CancelColor}{\color{red}}

\newcommand{\schwa}{ə}
\newcommand{\Schwa}{Ə}
\newcommand{\lette}[1]{\enquote*{#1}}
\newcommand{\virgo}[1]{\enquote*{#1}}
\newcommand{\definizione}[1]{\textcolor{red}{\textbf{#1}}}
\newcommand{\citaz}[2]{\guillemotleft#1\guillemotright\begin{flushright}#2\end{flushright}}
\newcommand{\cit}[1]{\guillemotleft#1\guillemotright}
\newcommand{\boh}[1]{\textcolor{red}{#1}}
\newcommand{\bianco}[1]{\textcolor{white}{#1}}

% ============================================================
%  PACCHETTI GRAFICI E VARI
% ============================================================
\usepackage{graphicx}
\usepackage{array}
\usepackage{booktabs}
\usepackage{placeins}
\usepackage{titling}
\usepackage{float}
\usepackage{wrapfig}
\usepackage{sidecap}
\usepackage{subfig}
\usepackage{subcaption}
\usepackage{floatflt}
\usepackage{animate}
\usepackage{standalone}
\usepackage{listings}
\usepackage{footmisc}
\usepackage{xstring}
\usepackage{ifthen}
\usepackage[strict]{changepage}
\usepackage[margin=1cm]{caption}

% ============================================================
%  HYPERREF (sempre per ultimo)
% ============================================================
\usepackage[colorlinks]{hyperref}
\hypersetup{linkcolor=magenta}
